\section{LLM评审者权重学习}

\subsection{问题建模与实验设计}

给定$N$个待评审样本，第$i$个样本有$K$个LLM评审者的打分$\mathbf{x}_i = (x_{i,1}, \ldots, x_{i,K})^\top$及对应的人工评审分数$y_i$。我们希望找到一组权重，使LLM评审的加权结果尽可能逼近人工评审。实验对比两种方案。

方案1采用凸组合权重，直接加权平均：$\hat{y}_i = \sum_{j=1}^{K} w_j x_{i,j}$，约束$w_j \geq 0$且$\sum_j w_j = 1$。通过最小化MSE求解：
\begin{equation}
\min_{\mathbf{w}} \frac{1}{N} \sum_{i=1}^{N} \left( \sum_{j=1}^{K} w_j x_{i,j} - y_i \right)^2
\end{equation}

方案2在加权后再做线性变换：
\begin{equation}
\hat{y}_i = a \left( \sum_{j=1}^{K} w_j x_{i,j} \right) + b
\end{equation}
权重$\mathbf{w}$保持相同约束，$(a, b)$通过最小二乘拟合。这样可以同时校正LLM分数的尺度和偏移。在三个评审维度（有效性、无干扰性、可部署性）上分别拟合独立权重。评审者打分从排名转为分数时采用组内归一化：$\text{score} = (N_{\text{max}} - \text{avg\_rank}) / (N_{\text{max}} - 1)$，其中$N_{\text{max}}$为该组的候选数上限。用留一法（LOO）和5折交叉验证评估泛化能力。

\subsection{实验结果}

\begin{table}[htbp]
\centering
\caption{两种方案的交叉验证结果}
\label{tab:cv-results}
\begin{tabular}{lccc}
\hline
\textbf{方案/维度} & \textbf{训练MSE} & \textbf{LOO MSE} & \textbf{5-fold MSE} \\
\hline
\multicolumn{4}{l}{\textit{凸组合权重}} \\
~~有效性   & 0.0344 & 0.0496 & 0.0458 \\
~~无干扰性 & 0.0344 & 0.0324 & 0.0329 \\
~~可部署性 & 0.0344 & 0.0341 & 0.0332 \\
~~\textbf{总体} & \textbf{0.0344} & \textbf{0.0387} & \textbf{0.0373} \\
\hline
\multicolumn{4}{l}{\textit{带校准权重}} \\
~~有效性   & 0.0263 & 0.0371 & 0.0335 \\
~~无干扰性 & 0.0263 & 0.0285 & 0.0256 \\
~~可部署性 & 0.0263 & 0.0378 & 0.0402 \\
~~\textbf{总体} & \textbf{0.0263} & \textbf{0.0345} & \textbf{0.0331} \\
\hline
改善幅度 & 23.6\% & 10.9\% & 11.3\% \\
\hline
\end{tabular}
\end{table}

从表\ref{tab:cv-results}可以看出，带校准的方案在5-fold验证中整体MSE为0.0331，比凸组合权重的0.0373降低了11\%。训练误差的大幅下降（23.6\%）与验证误差的温和下降说明校准参数$(a, b)$确实有助于拟合，但也需要警惕过拟合。分维度来看差异明显：有效性维度改善最大（5-fold: 0.0458→0.0335，降低27\%），这可能是因为LLM在判断漏洞利用效果时存在系统性的尺度偏差，需要线性变换来对齐人工标准；无干扰性维度也有22\%的改善；但可部署性维度反而变差（0.0332→0.0402），检查发现该维度下各LLM评审者分数高度一致（均匀权重MSE已达0.0339），此时引入校准参数反而带来过拟合风险。

\subsection{权重分布分析}

表\ref{tab:weight-distribution}给出了最终学到的权重（基于全量数据）。

\begin{table}[htbp]
\centering
\caption{学习到的评审者权重}
\label{tab:weight-distribution}
\begin{tabular}{lccc}
\hline
\textbf{评审者} & \textbf{有效性} & \textbf{无干扰性} & \textbf{可部署性} \\
\hline
ali-qwen  & 1.00 & 1.00 & 0.20 \\
deepseek  & 0.00 & 0.00 & 0.20 \\
glm4.5    & 0.00 & 0.00 & 0.20 \\
gpt5      & 0.00 & 0.00 & 0.20 \\
grok4     & 0.00 & 0.00 & 0.20 \\
\hline
\end{tabular}
\end{table}

有意思的是，有效性和无干扰性两个维度的权重完全集中在ali-qwen上，而可部署性呈均匀分布。这种"权重塌缩"现象在带校准的模型中其实合理：因为有$(a, b)$做线性变换，选择"最接近人工标准"的那个评审者再校准，往往比混合多个评审者更有效。可部署性的均匀分布则说明在这个维度上各评审者差异不大，或者它们与人工判断的相关性都较弱。

表\ref{tab:calibration-params}展示了校准参数。有效性和无干扰性的$a \approx 0.3\sim 0.5$，意味着LLM分数需要压缩一半左右才能对齐人工分数范围。可部署性的$a \approx 0$更极端——模型几乎放弃了LLM信号，直接预测人工评审的均值（$b = 0.5$）。

\begin{table}[htbp]
\centering
\caption{校准参数与拟合质量}
\label{tab:calibration-params}
\begin{tabular}{lcccc}
\hline
\textbf{维度} & $a$ & $b$ & \textbf{训练MSE} & \textbf{5-fold MSE} \\
\hline
有效性     & 0.50 & 0.25 & 0.025 & 0.034 \\
无干扰性   & 0.28 & 0.36 & 0.020 & 0.026 \\
可部署性   & 0.00 & 0.50 & 0.034 & 0.040 \\
\hline
\end{tabular}
\end{table}

\subsection{局限与改进空间}

当前方法的主要限制在于样本量偏小——每个维度只有20个数据点（5个LLM $\times$ 4种消融设置）。这直接导致可部署性维度几乎无法学到有效信号。权重塌缩到单个评审者虽然在数学上合理，但也反映出样本不足以支撑更精细的权重分配。现在只用了LLM的排名分数作为特征，如果能引入更多信号（比如评审置信度、多轮反馈中的分数变化、生成文本的长度和复杂度等），可能会提高拟合质量。非线性校准（分段函数、核方法等）也值得尝试，尤其是当LLM分数与人工分数的关系并非简单线性时。为了避免权重过度集中，可以加$L_2$正则项$\lambda \|\mathbf{w} - \mathbf{w}_{\text{uniform}}\|^2$。初步测试$\lambda = 0.1$时权重分布更均匀，但拟合效果略有下降。这需要在可解释性和准确性之间权衡。

通过交叉验证对比发现，在凸组合权重基础上增加线性校准能将MSE从0.037降到0.033（改善11\%）。有效性和无干扰性维度显著受益，但可部署性维度因LLM评审与人工判断相关性弱而效果不佳。权重倾向于集中在单个评审者上，这在小样本+线性校准的框架下是自然结果。后续可通过扩大数据集、引入正则化、尝试非线性模型来提升泛化能力和权重分布的稳健性。